\documentclass{article}

\usepackage[margin=1.2cm, top=1cm, bottom=1.8cm]{geometry}
\pagestyle{plain}
\linespread{1.4}

\usepackage{fontspec}
\usepackage{unicode-math}
\usepackage[main=greek, english]{babel}

\setmainfont{CMU Serif}
\setsansfont{CMU Sans Serif}
\setmonofont{CMU Typewriter Text}
\setmathfont{Latin Modern Math}

\usepackage{amsmath, enumitem}
\usepackage{dirtytalk, caption, subcaption, array, graphicx}
\usepackage{titlesec}
\usepackage{soul}
\usepackage[unicode]{hyperref}

\usepackage{pdfpages}

\hypersetup{
    hidelinks,
    bookmarksopen=true,
    bookmarksnumbered=true,
    pdftitle={Συστήματα Αναμονής - 1η Εργαστηριακή Άσκηση},
    pdfauthor={Πέππας Μιχαήλ-Αθανάσιος}
}

\newcolumntype{P}[1]{>{\centering\arraybackslash}p{#1}}
\titleformat*{\subsubsection}{\large\bfseries}

\begin{document}

\begin{center}
    {
    \Large
    {
    \textbf{ΕΘΝΙΚΟ ΜΕΤΣΟΒΙΟ ΠΟΛΥΤΕΧΝΕΙΟ\\
    ΣΧΟΛΗ ΗΛΕΚΤΡΟΛΟΓΩΝ ΜΗΧΑΝΙΚΩΝ\\
    ΚΑΙ ΜΗΧΑΝΙΚΩΝ ΥΠΟΛΟΓΙΣΤΩΝ\\}
    }

    \vspace{1.5cm}
    \rule{18cm}{0.5pt}

    Τομέας Επικοινωνιών, Ηλεκτρονικής και Συστημάτων Πληροφορικής\\
    Εργαστήριο Διαχείρισης και Βέλτιστου Σχεδιασμού Δικτύων Τηλεματικής - NETMODE\\

    \rule{18cm}{0.5pt}

    \textbf{Συστήματα Αναμονής}\\
    \textbf{$\mathbf{1^{η}}$ Εργαστηριακή Άσκηση}\\
    }
\end{center}

\vspace{0.8cm}
\begin{center}
\includegraphics[width=7cm]{EMP Logo.png}
\end{center}

\Large
{
\vspace{0.8cm}
\noindent
\textbf{Στοιχεία Φοιτητή}
\begin{itemize}
    \item Όνομα: Πέππας Μιχαήλ-Αθανάσιος
    \item Αριθμός Μητρώου: 03121026
    \item $3^{\eta}$ Εργαστηριακή Ομάδα (Δευτέρα, 12:45-14:30)
    \item Ημερομηνία Διεξαγωγής Εργαστηρίου: 11.03.2026
    \item Ημερομηνία Παράδοσης Αναφοράς: 16.03.2026
\end{itemize}
}

\newpage

\tableofcontents

\newpage

\section{Κατανομή Poisson}

Στο πρώτο μέρος της εργαστηριακής άσκησης μελετάται η κατανομή Poisson και ο τρόπος με τον
οποίο συνδέεται με τη διωνυμική κατανομή, τόσο θεωρητικά όσο και υπολογιστικά. Ειδικότερα,
ζητείται η σχεδίαση των συναρτήσεων μάζας πιθανότητας για διάφορες τιμές της παραμέτρου
\(\lambda\), η ανάδειξη της κατανομής Poisson ως ορίου της διωνυμικής, ο υπολογισμός της
μέσης τιμής και της διακύμανσής της, καθώς και η μελέτη της διάσπασης και υπέρθεσης κατανομών Poisson σε ένα απλό σενάριο τηλεφωνικού κέντρου.

\subsection{Συνάρτηση μάζας πιθανότητας της κατανομής Poisson}

Στο Διάγραμμα~\ref{fig:poisson_pmf} παρουσιάζονται οι γραφικές παραστάσεις των
συναρτήσεων μάζας πιθανότητας της κατανομής Poisson για παραμέτρους
\(\lambda = [7, 20, 40, 70]\), ως συνάρτηση του \(k\), για \(k = 0,1,2,\ldots,100\).

Αν μία διακριτή τυχαία μεταβλητή \(X\) ακολουθεί κατανομή Poisson με παράμετρο \(\lambda > 0\),
τότε η συνάρτηση μάζας πιθανότητας είναι:
\begin{equation}
    \Pr(X = k) = e^{-\lambda}\frac{\lambda^k}{k!},
    \qquad k = 0,1,2,\ldots
\end{equation}

\begin{figure}[htbp]
    \centering
    \includegraphics[width=0.92\textwidth]{part1_fig1.png}
    \caption{Συνάρτηση μάζας πιθανότητας της κατανομής Poisson για \(\lambda = 7, 20, 40, 70\).}
    \label{fig:poisson_pmf}
\end{figure}

\newpage
Από το παραπάνω διάγραμμα παρατηρούμε ότι όσο αυξάνεται η παράμετρος \(\lambda\), η
κατανομή μετατοπίζεται προς μεγαλύτερες τιμές του \(k\) και απλώνεται σε ευρύτερο εύρος
τιμών. Παράλληλα, η μέγιστη τιμή της συνάρτησης μάζας πιθανότητας μειώνεται, καθώς η
συνολική πιθανότητα παραμένει μοναδιαία και κατανέμεται σε περισσότερες τιμές. Πράγματι,
σε κάθε περίπτωση ισχύει η γνωστή συνθήκη κανονικοποίησης:
\begin{equation}
    \sum_{k=0}^{\infty} \Pr(X = k) = 1
\end{equation}

Επιπλέον, για τις συγκεκριμένες ακέραιες τιμές της παραμέτρου \(\lambda\), το μέγιστο της
κατανομής παρατηρείται σε δύο διαδοχικές θέσεις, στις \(k=\lambda-1\) και \(k=\lambda\),
γεγονός που συμφωνεί πλήρως με τη θεωρία της κατανομής Poisson. Άρα, όχι μόνο η θέση του
κέντρου της κατανομής μετακινείται προς τα δεξιά καθώς αυξάνει το \(\lambda\), αλλά και η
μορφή της γίνεται προοδευτικά πιο ομαλή.

\subsection{Κατανομή Poisson ως το όριο μιας διωνυμικής κατανομής}

Έστω ότι η τυχαία μεταβλητή \(X\) ακολουθεί διωνυμική κατανομή με παραμέτρους \(n\) και \(p\),
δηλαδή:
\begin{equation}
    \Pr(X = k) = \binom{n}{k} p^k (1-p)^{n-k},
    \qquad k = 0,1,\ldots,n
\end{equation}

Γνωρίζουμε ότι η κατανομή Poisson προκύπτει ως το όριο της διωνυμικής όταν:
\[
n \to \infty, \qquad p \to 0, \qquad np = \lambda \text{ σταθερό}
\]

Θέτοντας \(p = \lambda/n\), παίρνουμε:
\begin{equation}
    \Pr(X = k)
    =
    \binom{n}{k}
    \left(\frac{\lambda}{n}\right)^k
    \left(1-\frac{\lambda}{n}\right)^{n-k}
\end{equation}
και, λαμβάνοντας το όριο για \(n \to \infty\), προκύπτει:
\begin{equation}
    \Pr(X = k)
    \longrightarrow
    e^{-\lambda}\frac{\lambda^k}{k!}
\end{equation}
αφού:
\begin{equation}
    \lim_{n\to\infty}\left(1-\frac{\lambda}{n}\right)^n = e^{-\lambda}
\end{equation}

Συνεπώς, η κατανομή Poisson μπορεί να χρησιμοποιηθεί ως προσέγγιση της διωνυμικής όταν
το \(n\) είναι αρκετά μεγάλο και το \(p\) αρκετά μικρό, με σταθερό το γινόμενο \(np\).

Στην παρούσα άσκηση επιλέγουμε \(\lambda = 40\) και σχεδιάζουμε τη διωνυμική κατανομή για:
\[
n = [40, 120, 360, 1080, 40000]
\]
με αντίστοιχες πιθανότητες:
\[
p = \left[1,\ \frac{1}{3},\ \frac{1}{9},\ \frac{1}{27},\ \frac{1}{1000}\right]
\]
ώστε σε κάθε περίπτωση να ισχύει \(np = 40\).

\begin{figure}[htbp]
    \centering
    \includegraphics[width=0.92\textwidth]{part1_fig2.png}
    \caption{Η κατανομή Poisson ως όριο της διωνυμικής κατανομής για \(\lambda = 40\).}
    \label{fig:binomial_limit}
\end{figure}

Στο άνω υποδιάγραμμα του Διαγράμματος~\ref{fig:binomial_limit}, φαίνεται η μορφή των
διωνυμικών κατανομών για τα παραπάνω ζεύγη \((n,p)\). Η περίπτωση \(n=40\), \(p=1\) είναι
εκφυλισμένη, αφού η τυχαία μεταβλητή λαμβάνει πάντοτε την τιμή \(40\). Καθώς όμως το
\(n\) αυξάνεται και το \(p\) μειώνεται, με σταθερό το γινόμενο \(np=40\), η διωνυμική
κατανομή προσεγγίζει ολοένα και περισσότερο την αντίστοιχη Poisson.

Στο κάτω υποδιάγραμμα η σύγκλιση αυτή γίνεται ακόμη πιο εμφανής, αφού συγκρίνονται οι
κατανομές για \(n=120, 360, 1080, 40000\) με την πραγματική κατανομή \(\mathrm{Poisson}(40)\).
Παρατηρούμε ότι για πολύ μεγάλα \(n\) και πολύ μικρά \(p\), οι καμπύλες σχεδόν
ταυτίζονται. Άρα, το θεωρητικό οριακό αποτέλεσμα επιβεβαιώνεται και υπολογιστικά.

\subsection{Μέση τιμή και διακύμανση κατανομής Poisson}

Για την κατανομή Poisson με παράμετρο \(\lambda\), είναι γνωστό ότι:
\begin{equation}
    \mathbb{E}[X] = \lambda
    \qquad \text{και} \qquad
    \mathrm{Var}(X) = \lambda
\end{equation}

Στην παρούσα άσκηση επιλέγουμε την κατανομή \(\mathrm{Poisson}(40)\) και υπολογίζουμε, με
χρήση του Octave, τη μέση τιμή και τη διακύμανσή της. Η έξοδος του προγράμματος δίνεται
παρακάτω:

\begin{verbatim}
Mean value of Poisson with parameter lambda = 40 is
mean_value_pois = 40.000
Variance of Poisson with parameter lambda = 40 is
variance_pois = 40.000
\end{verbatim}

Άρα επιβεβαιώνεται πλήρως η θεωρία, αφού:
\[
\mathbb{E}[X] = 40,
\qquad
\mathrm{Var}(X) = 40
\]

Στη συνέχεια, υπολογίζουμε τη μέση τιμή και τη διακύμανση των αντίστοιχων διωνυμικών
κατανομών του προηγούμενου ερωτήματος. Η αριθμητική έξοδος είναι:

\begin{verbatim}
Mean value of the Binomial with parameter n and p
mean_value_bino =

   40   40   40   40   40

Variance of the Binomial as n increases
variance_bino =

         0   26.6667   35.5556   38.5185   39.9600
\end{verbatim}

Άρα, για κάθε μία από τις εξεταζόμενες διωνυμικές κατανομές, έχουμε:
\begin{itemize}
    \item Για \(n = 40\) και \(p = 1\): \(\mathbb{E}[X] = 40\), \(\mathrm{Var}(X) = 0\)
    \item Για \(n = 120\) και \(p = 1/3\): \(\mathbb{E}[X] = 40\), \(\mathrm{Var}(X) = 26.6667\)
    \item Για \(n = 360\) και \(p = 1/9\): \(\mathbb{E}[X] = 40\), \(\mathrm{Var}(X) = 35.5556\)
    \item Για \(n = 1080\) και \(p = 1/27\): \(\mathbb{E}[X] = 40\), \(\mathrm{Var}(X) = 38.5185\)
    \item Για \(n = 40000\) και \(p = 1/1000\): \(\mathbb{E}[X] = 40\), \(\mathrm{Var}(X) = 39.9600\)
\end{itemize}

Παρατηρούμε ότι η μέση τιμή της διωνυμικής κατανομής παραμένει σταθερά ίση με 40, αφού:
\[
\mathbb{E}[X_{\mathrm{bin}}] = np = 40
\]
σε όλες τις περιπτώσεις.

Για τη διακύμανση της διωνυμικής ισχύει:
\begin{equation}
    \mathrm{Var}(X_{\mathrm{bin}}) = np(1-p)
\end{equation}
οπότε, καθώς \(p \to 0\), προκύπτει:
\[
\mathrm{Var}(X_{\mathrm{bin}}) = np(1-p) \to np = \lambda
\]
και συνεπώς η διακύμανση της διωνυμικής τείνει στη διακύμανση της Poisson. Αυτό φαίνεται
και αριθμητικά, αφού οι τιμές \(0\), \(26.6667\), \(35.5556\), \(38.5185\), \(39.9600\)
προσεγγίζουν την τιμή \(40\) όσο το \(n\) αυξάνεται.

Συνεπώς, η σύγκλιση της διωνυμικής προς την Poisson δεν παρατηρείται μόνο γραφικά, αλλά
και μέσω των δύο πρώτων ροπών των αντίστοιχων κατανομών.

\subsection{Διάσπαση κατανομής Poisson - Υπέρθεση κατανομών}

Έστω ότι η τυχαία μεταβλητή \(X\) εκφράζει τον αριθμό κλήσεων σε ένα τηλεφωνικό κέντρο
στη διάρκεια μίας ώρας και ακολουθεί κατανομή Poisson με παράμετρο \(\lambda = 30\)
κλήσεις ανά ώρα. Κάθε κλήση είναι, ανεξάρτητα από τις υπόλοιπες, προς εξωτερικό με
πιθανότητα \(p = 0.3\) και προς εσωτερικό με πιθανότητα \(q = 1-p = 0.7\).

Αν \(X_1\) είναι ο αριθμός κλήσεων εξωτερικού και \(X_2\) ο αριθμός  κλήσεων εσωτερικού
ανά ώρα, τότε από την ιδιότητα διάσπασης της κατανομής Poisson προκύπτει:
\[
X_1 \sim \mathrm{Poisson}(p\lambda) = \mathrm{Poisson}(9)
\]
\[
X_2 \sim \mathrm{Poisson}(q\lambda) = \mathrm{Poisson}(21)
\]

Από την άλλη πλευρά, αφού:
\begin{equation}
    X = X_1 + X_2
\end{equation}
ισχύει:
\begin{equation}
    \mathbb{E}[X] = \mathbb{E}[X_1] + \mathbb{E}[X_2]
    \qquad \Longrightarrow \qquad
    \lambda = \lambda_1 + \lambda_2
\end{equation}
δηλαδή:
\[
30 = 9 + 21
\]

Το κοινό διάγραμμα των συναρτήσεων μάζας πιθανότητας των \(X_1\), \(X_2\) και \(X\)
φαίνεται στο Διάγραμμα~\ref{fig:poisson_split}.

\begin{figure}[htbp]
    \centering
    \includegraphics[width=0.92\textwidth]{part1_fig3.png}
    \caption{Κοινό διάγραμμα των \(\mathrm{Poisson}(9)\), \(\mathrm{Poisson}(21)\) και \(\mathrm{Poisson}(30)\).}
    \label{fig:poisson_split}
\end{figure}

Από το παραπάνω διάγραμμα παρατηρούμε ότι η κατανομή των κλήσεων εξωτερικού είναι
συγκεντρωμένη σε μικρότερες τιμές του \(k\), αφού έχει μικρότερη μέση τιμή
\((\lambda_1 = 9)\), ενώ η κατανομή των κλήσεων εσωτερικού μετατοπίζεται δεξιότερα, με
μέση τιμή \(\lambda_2 = 21\). Η συνολική κατανομή \(X\) έχει μέση τιμή \(\lambda = 30\)
και βρίσκεται ακόμη πιο δεξιά, όπως αναμένεται θεωρητικά, ως άθροισμα.

\paragraph{(β) Πιθανότητα οι επόμενες δύο κλήσεις να είναι εξωτερικού.}

Η πιθανότητα και οι δύο από τις επόμενες δύο κλήσεις να είναι εξωτερικού είναι:
\begin{equation}
    P_{\beta} = p^2 = (0.3)^2 = 0.09
\end{equation}

\paragraph{(γ) Πιθανότητα, δεδομένου ότι έγιναν ακριβώς 5 κλήσεις, οι 2 να ήταν εξωτερικού.}

Δεδομένου ότι υπήρξαν ακριβώς 5 κλήσεις, ο αριθμός των εξωτερικών κλήσεων ακολουθεί
διωνυμική κατανομή με παραμέτρους \(n=5\) και \(p=0.3\). Επομένως:
\begin{equation}
    P_{\gamma}
    =
    \binom{5}{2}(0.3)^2(0.7)^3
    =
    10 \cdot 0.09 \cdot 0.343
    =
    0.3087
\end{equation}

Η αριθμητική έξοδος του προγράμματος επιβεβαιώνει τα παραπάνω αποτελέσματα:

\begin{verbatim}
The probability that both of the next two calls are external is
Prob_b = 0.090000
Check with binopdf
Prob_b_bino = 0.090000
The probability that out of the 5 calls, two calls exactly are external, is
Prob_c = 0.3087
Check with binopdf
Prob_c_bino = 0.3087
\end{verbatim}

Συνεπώς, το ερώτημα της διάσπασης και υπέρθεσης της κατανομής Poisson επιβεβαιώνει τόσο
θεωρητικά όσο και υπολογιστικά ότι:
\begin{itemize}
    \item η διάσπαση μιας διαδικασίας \(\mathrm{Poisson}(\lambda)\) με πιθανότητες \(p\) και \(q\)
    οδηγεί σε δύο ανεξάρτητες κατανομές Poisson με παραμέτρους \(p\lambda\) και \(q\lambda\),
    \item η υπέρθεση των δύο αυτών κατανομών δίνει ξανά την αρχική κατανομή
    \(\mathrm{Poisson}(\lambda)\),
    \item οι επιμέρους πιθανότητες των εξωτερικών κλήσεων υπολογίζονται φυσικά με
    διωνυμικό μοντέλο, δεδομένου του συνολικού πλήθους κλήσεων.
\end{itemize}

\newpage
\section{Εκθετική κατανομή}

Στο δεύτερο μέρος της εργαστηριακής άσκησης μελετάται η εκθετική κατανομή, η οποία
χρησιμοποιείται ευρέως στη θεωρία ουρών για την περιγραφή χρόνων μεταξύ διαδοχικών
γεγονότων. Ειδικότερα, εξετάζονται οι συναρτήσεις πυκνότητας πιθανότητας τριών ανεξάρτητων
εκθετικών τυχαίων μεταβλητών, η κατανομή του ελαχίστου τους, καθώς και η ιδιότητα
απώλειας μνήμης της εκθετικής κατανομής.

\subsection{Συνάρτηση πυκνότητας πιθανότητας της εκθετικής κατανομής}

Έστω τρεις ανεξάρτητες συνεχείς τυχαίες μεταβλητές \(X_1\), \(X_2\), \(X_3\) που
ακολουθούν εκθετική κατανομή με μέσες τιμές:
\[
\frac{1}{\lambda_1} = 0.5,\qquad
\frac{1}{\lambda_2} = 1.5,\qquad
\frac{1}{\lambda_3} = 4
\]

Άρα οι αντίστοιχοι ρυθμοί είναι:
\[
\lambda_1 = 2,\qquad
\lambda_2 = \frac{2}{3},\qquad
\lambda_3 = \frac{1}{4}
\]

Η συνάρτηση πυκνότητας πιθανότητας της εκθετικής κατανομής με παράμετρο \(\lambda\) είναι:
\begin{equation}
    f_X(x) =
    \begin{cases}
        \lambda e^{-\lambda x}, & x \geq 0 \\
        0, & x < 0
    \end{cases}
\end{equation}

Στο Διάγραμμα~\ref{fig:exp_pdf} παρουσιάζονται οι συναρτήσεις πυκνότητας πιθανότητας των
τριών παραπάνω τυχαίων μεταβλητών για τιμές του οριζόντιου άξονα από \(0\) έως \(8\).

\begin{figure}[htbp]
    \centering
    \includegraphics[width=0.92\textwidth]{part2_fig1.png}
    \caption{Συνάρτηση πυκνότητας πιθανότητας της εκθετικής κατανομής για τρεις διαφορετικές μέσες τιμές.}
    \label{fig:exp_pdf}
\end{figure}

Από το διάγραμμα παρατηρούμε ότι όσο μικρότερη είναι η μέση τιμή \(1/\lambda\), τόσο
μεγαλύτερος είναι ο ρυθμός \(\lambda\) και τόσο πιο απότομα φθίνει η πυκνότητα. Έτσι, η
τυχαία μεταβλητή με μέση τιμή \(0.5\) sec παρουσιάζει τη μεγαλύτερη αρχική τιμή στο
\(x=0\) και τη γρηγορότερη μείωση. Αντίθετα, η τυχαία μεταβλητή με μέση τιμή \(4\) sec
φθίνει πιο αργά και διατηρεί σημαντική πιθανότητα και για μεγαλύτερες τιμές του χρόνου.

Επομένως, η παράμετρος \(\lambda\) ελέγχει άμεσα το πόσο γρήγορα λαμβάνονται μικρές
τιμές από την εκθετική τυχαία μεταβλητή: μεγάλος ρυθμός σημαίνει μικρότερη μέση τιμή
και μεγαλύτερη συγκέντρωση πυκνότητας κοντά στο μηδέν.

\subsection{Συνδυασμός εκθετικών κατανομών}

Θεωρούμε τώρα την τυχαία μεταβλητή:
\[
X = \min(X_1, X_2, X_3)
\]
όπου οι \(X_1\), \(X_2\), \(X_3\) είναι ανεξάρτητες εκθετικές τυχαίες μεταβλητές όπως
παραπάνω.

Για να προσδιορίσουμε τη συνάρτηση κατανομής της \(X\), γράφουμε:
\begin{align}
    F_X(x)
    &= \Pr(X \leq x) \\
    &= 1 - \Pr(X > x) \\
    &= 1 - \Pr(X_1 > x,\ X_2 > x,\ X_3 > x)
\end{align}

Λόγω ανεξαρτησίας:
\begin{equation}
    \Pr(X_1 > x,\ X_2 > x,\ X_3 > x)
    =
    \Pr(X_1 > x)\Pr(X_2 > x)\Pr(X_3 > x)
\end{equation}

Για εκθετικές τυχαίες μεταβλητές ισχύει:
\[
\Pr(X_i > x) = e^{-\lambda_i x}, \qquad x \geq 0
\]
άρα:
\begin{align}
    F_X(x)
    &= 1 - e^{-\lambda_1 x} e^{-\lambda_2 x} e^{-\lambda_3 x} \\
    &= 1 - e^{-(\lambda_1+\lambda_2+\lambda_3)x}
\end{align}

Συνεπώς, το ελάχιστο τριών ανεξάρτητων εκθετικών τυχαίων μεταβλητών είναι επίσης εκθετικό:
\begin{equation}
    X = \min(X_1, X_2, X_3) \sim \mathrm{Exp}(\lambda_1+\lambda_2+\lambda_3)
\end{equation}

Στη συγκεκριμένη άσκηση:
\[
\lambda_1+\lambda_2+\lambda_3
=
2 + \frac{2}{3} + \frac{1}{4}
=
2.916667
\]
οπότε:
\[
X \sim \mathrm{Exp}(2.916667)
\]
και η μέση τιμή του είναι:
\begin{equation}
    \mathbb{E}[X] = \frac{1}{\lambda_1+\lambda_2+\lambda_3}
    = \frac{1}{2.916667}
    = 0.342857\ \text{sec}
\end{equation}

Το αποτέλεσμα αυτό επιβεβαιώνεται και από την έξοδο του προγράμματος:

\begin{verbatim}
Rate of X = min(X1, X2, X3): lambda = 2.916667
Mean value of X = min(X1, X2, X3): 0.342857 sec
\end{verbatim}

Στο Διάγραμμα~\ref{fig:exp_cdf} παρουσιάζονται οι αθροιστικές συναρτήσεις κατανομής των
\(X_1\), \(X_2\), \(X_3\) και του \(X=\min(X_1,X_2,X_3)\).

\begin{figure}[htbp]
    \centering
    \includegraphics[width=0.92\textwidth]{part2_fig2.png}
    \caption{Αθροιστικές συναρτήσεις κατανομής των \(X_1\), \(X_2\), \(X_3\) και του \(X=\min(X_1,X_2,X_3)\).}
    \label{fig:exp_cdf}
\end{figure}

Από το παρακάτω διάγραμμα παρατηρούμε ότι η συνάρτηση κατανομής του ελαχίστου αυξάνεται
ταχύτερα από τις αντίστοιχες των αρχικών τυχαίων μεταβλητών. Αυτό σημαίνει ότι το
\(\min(X_1,X_2,X_3)\) λαμβάνει, κατά κανόνα, μικρότερες τιμές από καθεμία εκ των
\(X_1\), \(X_2\), \(X_3\) ξεχωριστά. Επομένως, η μέση τιμή του είναι μικρότερη από κάθε
μία από τις αρχικές μέσες τιμές \(0.5\), \(1.5\) και \(4\) sec, γεγονός που συμφωνεί
πλήρως με τη θεωρία.

\newpage
\subsection{Απώλεια μνήμης της εκθετικής κατανομής}

Η εκθετική κατανομή διαθέτει την πολύ σημαντική ιδιότητα απώλειας μνήμης, σύμφωνα με
την οποία:
\begin{equation}
    \Pr(X > t+s \mid X > s) = \Pr(X > t)
\end{equation}
για κάθε \(t,s \geq 0\).

Στην παρούσα άσκηση θεωρούμε μία συνεχόμενη τυχαία μεταβλητή \(X\) που ακολουθεί
εκθετική κατανομή με μέση τιμή:
\[
\frac{1}{\lambda} = 2.5\ \text{sec}
\]
και υπολογίζουμε τις πιθανότητες:
\[
\Pr(X > 20000)
\qquad \text{και} \qquad
\Pr(X > 45000 \mid X > 25000)
\]
όπου οι αριθμοί \(20000\), \(25000\), \(45000\) δηλώνουν θέσεις στο διάνυσμα
\[
k = 0:0.00001:8
\]
και όχι πραγματικούς χρόνους σε sec.

Από την έξοδο του προγράμματος έχουμε:

\begin{verbatim}
For X ~ Exp with mean = 2.5 sec:

Index 20000 corresponds to time t = 0.19999 sec
Index 25000 corresponds to time s = 0.24999 sec
Index 45000 corresponds to time t+s = 0.44999 sec

Pr(X > 20000) = 0.923120
Pr(X > 45000 | X > 25000) = 0.923116
Memoryless check Pr(X > (45000-25000)) = 0.923116
\end{verbatim}

Παρατηρούμε ότι οι δύο πιθανότητες είναι πρακτικά ίσες:
\[
\Pr(X > 20000) \approx 0.923120
\]
\[
\Pr(X > 45000 \mid X > 25000) \approx 0.923116
\]

Η πολύ μικρή διαφορά στα τελευταία δεκαδικά ψηφία οφείλεται αποκλειστικά στη
διακριτοποίηση του άξονα και στη στρογγυλοποίηση των αριθμητικών υπολογισμών.
Θεωρητικά, οι δύο πιθανότητες πρέπει να ταυτίζονται ακριβώς λόγω της ιδιότητας
απώλειας μνήμης.

Η ερμηνεία της ιδιότητας αυτής είναι ιδιαίτερα σημαντική στη θεωρία ουρών. Αν ένας
χρόνος αναμονής ή εξυπηρέτησης ακολουθεί εκθετική κατανομή, τότε το γεγονός ότι έχει
ήδη περάσει κάποιο χρονικό διάστημα χωρίς να συμβεί το γεγονός δεν αλλάζει την
κατανομή του υπολειπόμενου χρόνου. Με άλλα λόγια, το σύστημα δεν θυμάται πόσος χρόνος
έχει ήδη περάσει.

Συνεπώς, το συγκεκριμένο ερώτημα επιβεβαιώνει υπολογιστικά μία από τις βασικότερες
ιδιότητες της εκθετικής κατανομής, η οποία εξηγεί γιατί η εκθετική κατανομή
εμφανίζεται τόσο συχνά σε στοχαστικά μοντέλα άφιξης και εξυπηρέτησης.

\newpage
\section{Διαδικασία καταμέτρησης Poisson}

Στο τρίτο μέρος της εργαστηριακής άσκησης μελετάται η διαδικασία καταμέτρησης Poisson
\(N(t)\), δηλαδή η στοχαστική διαδικασία που περιγράφει το πλήθος γεγονότων που έχουν
συμβεί μέχρι τη χρονική στιγμή \(t\). Ειδικότερα, εξετάζονται οι χρόνοι μεταξύ
διαδοχικών γεγονότων, η μορφή της ίδιας της διαδικασίας καταμέτρησης, η εκτίμηση του
ρυθμού άφιξης \(\lambda\) από προσομοιωμένα δεδομένα, καθώς και η θεωρητική κατανομή του
πλήθους των γεγονότων σε δεδομένα χρονικά παράθυρα.

\subsection{Διαδικασία καταμέτρησης Poisson}

Έστω διαδικασία Poisson με ρυθμό:
\[
\lambda = 10\ \text{γεγονότα/sec}
\]

Από τη θεωρία είναι γνωστό ότι οι χρόνοι που μεσολαβούν ανάμεσα στην εμφάνιση δύο
διαδοχικών γεγονότων μιας διαδικασίας Poisson είναι ανεξάρτητες και ισόνομες
εκθετικές τυχαίες μεταβλητές με μέση τιμή:
\begin{equation}
    \mathbb{E}[T] = \frac{1}{\lambda}
\end{equation}

Στη συγκεκριμένη περίπτωση:
\[
\mathbb{E}[T] = \frac{1}{10} = 0.1\ \text{sec}
\]

Το παραπάνω επιβεβαιώνεται και από την έξοδο του προγράμματος:

\begin{verbatim}
For a Poisson process with rate lambda = 10.00 events/sec,
the interarrival times follow an exponential distribution
with mean 1/lambda = 0.1000 sec.
\end{verbatim}

Με χρήση της εντολής \texttt{exprnd()} στο Octave δημιουργήθηκαν 100 διαδοχικοί χρόνοι
μεταξύ αφίξεων και, στη συνέχεια, μέσω αθροιστικής πρόσθεσης προέκυψαν οι χρονικές
στιγμές εμφάνισης των γεγονότων. Από αυτές κατασκευάστηκε, με τη συνάρτηση
\texttt{stairs}, η διαδικασία καταμέτρησης Poisson \(N(t)\), η οποία φαίνεται στο
Διάγραμμα~\ref{fig:poisson_counting}.

\begin{figure}[htbp]
    \centering
    \includegraphics[width=0.92\textwidth]{part3_fig1.png}
    \caption{Διαδικασία καταμέτρησης Poisson για \(\lambda = 10\) γεγονότα/sec.}
    \label{fig:poisson_counting}
\end{figure}

Παρατηρούμε ότι η \(N(t)\) είναι μία κλιμακωτή, μη φθίνουσα συνάρτηση, η οποία
αυξάνεται κατά μία μονάδα κάθε φορά που συμβαίνει ένα νέο γεγονός. Τα οριζόντια
τμήματα του διαγράμματος αντιστοιχούν σε χρονικά διαστήματα στα οποία δεν συμβαίνει
καμία άφιξη, ενώ τα κατακόρυφα άλματα αντιστοιχούν στις χρονικές στιγμές αφίξεων.

Η συνολική μορφή της καμπύλης είναι πλήρως αναμενόμενη για μία διαδικασία Poisson. Αν
και η άνοδος δεν είναι απολύτως ομαλή, καθώς πρόκειται για τυχαία πραγματοποίηση,
παρατηρείται ότι όσο αυξάνεται ο χρόνος, το πλήθος των γεγονότων αυξάνει περίπου
γραμμικά, με μέσο ρυθμό κοντά στο \(\lambda = 10\) γεγονότα ανά δευτερόλεπτο.

\subsection{Μέσος αριθμός γεγονότων διαδικασίας Poisson}

Στη συνέχεια, εκτιμούμε τον μέσο ρυθμό άφιξης των γεγονότων μέσω του λόγου:
\begin{equation}
    \hat{\lambda} = \frac{n}{T_n}
\end{equation}
όπου \(n\) είναι το συνολικό πλήθος γεγονότων και \(T_n\) ο χρόνος εμφάνισης του
\(n\)-οστού γεγονότος.

Ο υπολογισμός αυτός πραγματοποιήθηκε αρχικά για \(n=100\) και έπειτα επαναλήφθηκε για:
\[
n = 200,\ 400,\ 700,\ 1000,\ 5000,\ 50000,\ 100000
\]

Η έξοδος του προγράμματος είναι η εξής:

\begin{verbatim}
Estimated rate lambda_hat for increasing number of events:
--------------------------------------------------------
n =    100   total_time =     9.410809 sec   lambda_hat =  10.626079 events/sec
n =    200   total_time =    19.676798 sec   lambda_hat =  10.164255 events/sec
n =    400   total_time =    40.251597 sec   lambda_hat =   9.937494 events/sec
n =    700   total_time =    70.536601 sec   lambda_hat =   9.923926 events/sec
n =   1000   total_time =   105.505041 sec   lambda_hat =   9.478220 events/sec
n =   5000   total_time =   490.597130 sec   lambda_hat =  10.191662 events/sec
n =  50000   total_time =  4999.763849 sec   lambda_hat =  10.000472 events/sec
n = 100000   total_time = 10033.501710 sec   lambda_hat =   9.966610 events/sec
\end{verbatim}

Τα παραπάνω αποτελέσματα αποτυπώνονται γραφικά στο Διάγραμμα~\ref{fig:lambda_estimation}.

\begin{figure}[htbp]
    \centering
    \includegraphics[width=0.92\textwidth]{part3_fig2.png}
    \caption{Εκτίμηση της παραμέτρου \(\lambda\) για αυξανόμενο πλήθος γεγονότων.}
    \label{fig:lambda_estimation}
\end{figure}

Για \(n=100\) προκύπτει η εκτίμηση:
\[
\hat{\lambda} = 10.626079
\]
η οποία αποτελεί μία καλή, αλλά όχι τέλεια, προσέγγιση της πραγματικής τιμής
\(\lambda = 10\). Αυτό είναι αναμενόμενο, επειδή για μικρό σχετικά δείγμα η
προσομοίωση επηρεάζεται περισσότερο από τη στοχαστική μεταβλητότητα των
μεσοδιαστημάτων άφιξης.

Καθώς όμως το πλήθος των γεγονότων αυξάνεται σημαντικά, η εκτίμηση
\(\hat{\lambda}\) σταθεροποιείται όλο και περισσότερο κοντά στην πραγματική τιμή
\(\lambda = 10\). Για παράδειγμα, για \(n=50000\) προκύπτει:
\[
\hat{\lambda} = 10.000472
\]
ενώ για \(n=100000\) προκύπτει:
\[
\hat{\lambda} = 9.966610
\]

Το φαινόμενο αυτό είναι πλήρως συμβατό με τον Νόμο των Μεγάλων Αριθμών: όσο
αυξάνεται το μέγεθος του δείγματος, τόσο η εμπειρική εκτίμηση πλησιάζει τη
θεωρητική τιμή της παραμέτρου.

Επομένως, από την προσομοίωση επιβεβαιώνεται ότι ο ρυθμός άφιξης της διαδικασίας
Poisson μπορεί να εκτιμηθεί με καλή ακρίβεια από τον λόγο πλήθους γεγονότων προς
συνολικό χρόνο παρατήρησης, και ότι η ακρίβεια αυτή βελτιώνεται όσο αυξάνεται ο
αριθμός των γεγονότων.

\subsection{Πλήθος γεγονότων σε χρονικό παράθυρο}

Από τη θεωρία της διαδικασίας Poisson γνωρίζουμε ότι το πλήθος των γεγονότων που
συμβαίνουν σε ένα χρονικό παράθυρο διάρκειας \(\Delta T\) ακολουθεί κατανομή
Poisson με παράμετρο:
\begin{equation}
    \lambda \Delta T
\end{equation}

Άρα, αν
\[
N(\Delta T) = \text{πλήθος γεγονότων σε χρονικό διάστημα διάρκειας } \Delta T
\]
τότε:
\begin{equation}
    N(\Delta T) \sim \mathrm{Poisson}(\lambda \Delta T)
\end{equation}
και συνεπώς:
\begin{equation}
    \mathbb{E}[N(\Delta T)] = \lambda \Delta T
\end{equation}

Στην παρούσα άσκηση θεωρούμε δύο μη επικαλυπτόμενα χρονικά παράθυρα:
\[
\Delta T_1 = 5\ \text{sec}
\qquad \text{και} \qquad
\Delta T_2 = 10\ \text{sec}
\]

Για το πρώτο παράθυρο:
\[
N_1 \sim \mathrm{Poisson}(\lambda \Delta T_1)
= \mathrm{Poisson}(10 \cdot 5)
= \mathrm{Poisson}(50)
\]
άρα:
\[
\mathbb{E}[N_1] = 50
\]

Για το δεύτερο παράθυρο:
\[
N_2 \sim \mathrm{Poisson}(\lambda \Delta T_2)
= \mathrm{Poisson}(10 \cdot 10)
= \mathrm{Poisson}(100)
\]
άρα:
\[
\mathbb{E}[N_2] = 100
\]

Επειδή τα δύο χρονικά διαστήματα είναι μη επικαλυπτόμενα, το συνολικό πλήθος
γεγονότων στα δύο παράθυρα είναι:
\[
N_1 + N_2
\]
και η μέση τιμή του είναι:
\begin{equation}
    \mathbb{E}[N_1 + N_2]
    =
    \mathbb{E}[N_1] + \mathbb{E}[N_2]
    =
    50 + 100
    =
    150
\end{equation}

Η έξοδος του προγράμματος επιβεβαιώνει ακριβώς τα παραπάνω:

\begin{verbatim}
Theoretical results for two non-overlapping time windows:
--------------------------------------------------------
Window 1: DeltaT1 = 5 sec  -> mean number of events = 50.000000
Window 2: DeltaT2 = 10 sec  -> mean number of events = 100.000000
Total over both windows      -> mean number of events = 150.000000

Distribution remarks:
N1 ~ Poisson(lambda * DeltaT1) = Poisson(50.00)
N2 ~ Poisson(lambda * DeltaT2) = Poisson(100.00)
N1 + N2 ~ Poisson(150.00)
\end{verbatim}

Συνεπώς, το πλήθος των γεγονότων σε χρονικό παράθυρο μιας διαδικασίας Poisson
ακολουθεί κατανομή Poisson με παράμετρο ίση με τον ρυθμό επί τη διάρκεια του
παραθύρου, ενώ για μη επικαλυπτόμενα χρονικά διαστήματα οι μέσοι αριθμοί γεγονότων
αθροίζονται. Στη συγκεκριμένη περίπτωση, ο θεωρητικός συνολικός μέσος αριθμός
γεγονότων στα δύο χρονικά παράθυρα είναι ίσος με \(150\).

\newpage
\section{Παράρτημα Υλοποιήσεων σε Octave}

\subsection{Matlab Code - Part 1}
\includepdf[pages=-]{lab1_part1.pdf}

\subsection{Matlab Code - Part 2}
\includepdf[pages=-]{lab1_part2.pdf}

\subsection{Matlab Code - Part 3}
\includepdf[pages=-]{lab1_part3.pdf}

\end{document}
